\section{Introduction}
Multi-agent systems, as upgraded extensions of single-agent systems, require collaborative efforts from multiple agents working in concert~\citep{agent-survey-1, agent-survey-2}. 
With the advancement of Large Language Models (LLMs)~\citep{instructGPT, gpt-4, llama, llama2}, multi-agent applications have made great progress in both research and industrial communities, including software engineering~\citep{meta-gpt}, society simulation~\citep{aitown}, and intelligent assistant~\citep{auto-gen, autogpt}.
Although significant progress has been made in multi-agent scenarios, there are still major challenges remaining in multi-agent application development.  

\textit{Developing a multi-agent application is more complex than creating a single-agent one.}
Unlike single-agent setups where an agent solely interacts with users, the development in the multi-agent scenario requires careful creation and management of multiple models and agents~\citep{agent-survey-1, agent-survey-2}, which poses high requirements for both versatility and handiness for a platform.
In particular, the following aspects feature the challenges: 1) Agents involved in a multi-agent application can specialize at different functions via different initial configurations; 2) A multi-agent application may require agents to be executed in a standardized operating procedure (SOP) or a more dynamic workflow; 3)The communication pattern between agents can be varying from one-to-one or broadcasting (e.g., a discussion group of agents).
As a result, developers expect a handy platform that can provide concise and clear programming patterns when taking care of all the aspects above, accelerating and facilitating the development cycle.
Achieving versatility and handiness simultaneously requires careful design and taking trade-offs, and it remains a persistent goal for all multi-agent platform designs.

\textit{Aberrations are tinderboxs in a multi-agent system.}
Although LLMs have advanced rapidly, they still struggle with issues like hallucination~\citep{hall1, hall2} and inadequate instruction-following~\citep{follow1, follow2}.
Besides, an agent can be equipped with various tools, but those tools introduce additional uncertainties (e.g., accessibility to a database or the search engine).
From the perspective of multi-agent system robustness, any unexpected error or response can propagate to the whole system, causing a series of cascading effects if not handled properly.
Thus, it is crucial for multi-agent applications to autonomously detect and handle unexpected responses from LLMs. 
While LLMs may assist in identifying and managing these errors, it remains a challenge to determine whether they can resolve errors on their own and to automatically provide the necessary information for error correction. 
Consequently, designing fault-tolerant that incorporate LLMs is a key challenge in the development of multi-agent applications. 

\textit{Supporting agents with multi-modal data, tools, and external knowledge is highly systematic.} 
Besides generating answers with LLMs, agents are expected to be more versatile, including generating and handling multi-modal data~\citep{multimodal, dalle3}, preparing and invoking functions as tools~\citep{react, shen2024small}, managing external knowledge banks, and using the retrieved knowledge for augmentation generation~\citep{lewis2020retrieval}.
However, integrating these functionalities in multi-agent applications requires a comprehensive and systematic approach. 
Supporting multi-modal content is a complex endeavor, necessitating considerations for data storage, presentation, user interaction, message transmission, and communication. 
Tool utilization of agents requires unifying the function calling pattern and output parsing, prompting to instruct LLMs, and designing reasoning mechanisms to ensure the tasks can be accomplished step by step.
As for external knowledge, beyond the retrieval-augmented generation (RAG) techniques, we need to consider how to efficiently share and manage the knowledge in multi-agent scenarios while leaving enough flexibility for retrieval strategies.
While some existing works investigate how those techniques individually work within specialized agent systems, general platform-level programming interfaces remain absent.
% There is an increasing number of agent or LLM applications targeting multi-modal content generation. 
% Supporting multi-modal data~\citep{multimodal, dalle3} in multi-agent applications requires a comprehensive and systematic approach, which takes considerations for data storage, presentation, user interaction, message transmission, communication, \etc
% However, meeting all these requirements presents some new challenges, including ensuring data consistency across different formats, maintaining high performance during data transmission and agent communication, and avoiding the introduction of complex concepts for both developers and users. 
% While recent research has introduced some solutions tailored to specific applications, general platform-level programming interfaces to promote multi-modal applications remain absent.

\textit{Distributed applications bring extra programming difficulties and system design challenges.}
An industrial-oriented scenario for multi-agent applications is that the agents are owned by different organizations and run on different machines because the agents are equipped with unique private knowledge or patented tools.
Developing such applications usually requires the developers to have professional knowledge of distributed system programming and optimization in the design phase.
Besides, distributed applications usually require a great extra effort in the development and testing, especially when debugging and diagnosing issues spread across distributed processes or agents.
Moreover, integrating advanced features like multi-modal data processing poses additional challenges in a distributed setting, when the agents require different time to accomplish the sub-tasks or the generated contents are very heterogeneous. 
Poor distributed system design can result in excessive communication overhead between agents. 
% Therefore, it's critical but challenging for developers to address these issues effectively to ensure effective system operations.
Therefore, building distributed multi-agent applications requires the large efforts of experienced developers and a high barrier for beginners to migrate their prototypes to a distributed style for optimal efficiency.

To tackle the aforementioned challenges, we introduce \textbf{\ours}, a novel multi-agent platform designed for developers with varying levels of expertise. 
\ours is well-designed with a message exchange communication mechanism that embodies great usability, robustness, and efficiency.
We underscore the salient features of \ours as follows:

\textbf{Exceptional Usability for Developers}.
\ours is designed with a fundamental emphasis on ease of use, particularly for developers with varying levels of expertise. 
By implementing a procedure-oriented message exchange mechanism, \ours ensures a smooth learning curve on multi-agent application development. 
To alleviate the programming burdens, \ours offers an extensive suite of syntactic utilities, including various pipelines and an information-sharing mechanism.
Besides programming with our framework, we also improve usability by providing a \emph{zero-code drag-and-drop programming workstation}, which can enable those with limited Python programming experience to build their own applications with little effort.
Compared with building the skeleton of the application, prompt tuning can be a more time-consuming stage in multi-agent application development.
In \ours, we equip our agents with a set of \emph{automatic prompt tuning} mechanisms to relieve such burden.
Coupled with rich built-in resources and integrated user interaction modules, \ours makes building a multi-agent application much more enjoyable than ever. 

\textbf{Robust Fault Tolerance for Diverse LLMs and APIs}.
As the scale and scope of models and APIs expand, a robust fault-tolerance mechanism in multi-agent applications becomes paramount. 
\ours integrates a comprehensive service-level retry mechanism to maintain API reliability. 
\ours is equipped with a set of rule-based correction tools to handle some obvious formatting problems in the responses of LLMs. 
Moreover, \ours offers customizable fault tolerance configurations, enabling developers to tailor their own fault tolerance mechanism through parameters like \texttt{parse\_func}, \texttt{fault\_handler}, and \texttt{max\_retries}.
While admittedly, not all the errors can be handled by the aforementioned mechanism, we propose a logging system with customized features for multi-agent applications as the last safeguard for \ours. 

\textbf{Extensive Compatibility for Multi-Modal, Tools, and External Knowledge}. 
With the remarkable progress of large-scale multi-modal models, \ours supports multi-modal data (\eg texts, images, audio, and videos) in dialog conversation, message transmission, and data storage. 
Specifically, \ours decouples multi-modal data transmission from storage and employs \emph{a lazy loading strategy} by providing a unified URL-based attribute in messages. 
During message transmission, \ours only attaches a URL to the message, and the multi-modal data is loaded only when necessary, such as when being rendered in web UI or invoked by model wrappers. 
For tool usage, \ours provides a component, called \emph{service toolkit}, as a one-step solution for tool usage, including function preprocessing, prompt engineering, reasoning, and response parsing with fault-tolerance features.
To support efficient external knowledge usage, \ours provides end-to-end, highly configurable, and sharable knowledge processing modules for retrieval-augmented generation (RAG), from data preprocessing to customizable retrieval.

\textbf{Optimized Efficiency for Distributed Multi-Agent Operations}.
Recognizing the vital importance of distributed deployment, \ours introduces an actor-based distributed mechanism that enables centralized programming of complex distributed workflows, and automatic parallel optimization. 
Particularly, the workflows for local and distributed deployments is a exactly the same one, indicating negligible overhead when migrating applications between centralized and distributed environments. 
With such a distribution framework, \ours empowers developers to concentrate on the application design rather than implementation details. 

\noindent
\paragraph{Summary}
To summarize, \ours, a novel multi-agent platform proposed for flexibility and robustness, includes the following advanced features:
\begin{enumerate}
    \item 
    \ours provides a procedure-oriented message exchange mechanism with a set of syntactic features to facilitate multi-agent programming, a zero-code drag-and-drop programming workstation, and a set of automatic prompt tuning mechanisms.
    
    \item 
    The fault tolerance designs of \ours enable developers to handle errors elegantly for their applications.
    
    \item 
    The support for the multi-modal applications reduces the overheads of heterogeneous data generation and transmission. The service toolkit component facilitates the tool usage of agents in \ours, and the knowledge processing modules provide a flexible solution for agents to handle different information.
    
    \item
    The actor-based distributed mode of \ours can help develop efficient and reliable distributed multi-agent applications seamlessly.
\end{enumerate}

\noindent\paragraph{Roadmap}
In the following sections, we navigate through the core components and capabilities of \ours, showcasing its role in advancing the development and deployment of multi-agent applications. Section \ref{sec:overview} provides an overview, while Section \ref{sec:usability} focuses on the user experience.
Section \ref{sec:tolerance} introduces the fault tolerance mechanism in \ours.
Sections \ref{sec:mutil-modal}, \ref{sec:tools}, and \ref{sec:rag} cover the multi-modal support, tool usage, and retrieval-augmented generation modules in \ours. 
Section \ref{sec:distribute} presents our platform's support for distributed multi-agent applications.
Use cases are presented in Section \ref{sec:app}, related work is summarized in Section \ref{sec:related_work}, and concluding thoughts are recorded in Section \ref{sec:conclusion}.
